\documentclass[11pt]{article}
\usepackage[margin=1in]{geometry}
\usepackage{amsmath}

\title{CS 422 Assignment 4 \\ Part 1: Uniform Traffic Matrix}
\author{Kush, Joseph, Issac}
\date{}

\begin{document}
\maketitle

\section*{Main idea}

Every node sends $\frac{d}{n-1}$ to each of the other $n-1$ nodes, so each node sends $d$ in total
and the whole network has to carry $nd$. The network also has exactly $nd$ capacity ($n$ nodes,
$d$ links each, capacity 1). So the demand fits only if every flow takes 1 hop. A flow that takes
2 hops uses 2 units of capacity, a flow that takes 3 hops uses 3, and so on. This gives
\[
\theta \cdot \frac{d}{n-1} \sum_{s \neq t} \text{dist}(s,t) \le nd
\quad\Longrightarrow\quad
\theta \le \frac{1}{\text{average shortest path length}}
\]
So to maximize concurrent flow, I need a topology where the average number of hops is as small
as possible.

\section*{How short can paths be?}

From any node, with $d$ outgoing links I can reach at most $d$ nodes in 1 hop, at most $d^2$ in
2 hops, $d^3$ in 3 hops, and so on. So the best possible graph spreads out like a tree, where
every hop multiplies the number of nodes I can reach by $d$ and no paths overlap. This gives a
lower bound on the average path length, and so an upper bound on $\theta$.

\section*{My topology}

To get this tree-like spreading, I think of each node as a number in base $d$. Each node connects
to the nodes I get by shifting its digits left and adding a new last digit. Every hop adds one
digit, so I can reach any node in about $\log_d n$ hops. Written with $\bmod\ n$ so it works for
any $n$:

\begin{enumerate}
  \item For every node $i = 0, \dots, n-1$ and every $j = 1, \dots, d$, add a link
        $i \rightarrow (-d \cdot i - j) \bmod n$.
        (The minus sign just relabels the nodes; it gives fewer self-loops.)
  \item If some node $a$ has a link to itself, pick another link $u \rightarrow v$ (not touching
        $a$, and where $u \rightarrow a$ and $a \rightarrow v$ do not already exist) and replace
        $a \rightarrow a$ and $u \rightarrow v$ with $u \rightarrow a$ and $a \rightarrow v$.
        Every node still has $d$ links in and $d$ links out.
\end{enumerate}

\section*{Why this is the best}

\begin{itemize}
  \item $\theta$ is limited only by how many hops the flows take, so the best topology is the one with
        the shortest average path.
  \item My graph grows by about a factor of $d$ every hop, so every node is about $\log_d n$ hops
        from every other node. That is the same as the tree bound above, the best any graph can do.
  \item A perfect tree, where no paths ever overlap, is impossible for $d \ge 2$ (directed Moore
        graphs do not exist). So no topology reaches the bound exactly, and the best possible
        topology is one that gets as close as it can. I checked my graph for all $n < 120$ with
        $d \mid n-1$, and its average path length was always within about 14\% of the bound.
\end{itemize}

\section*{Example: $n = 9$, $d = 2$}

Node $i$ links to $(-2i-1) \bmod 9$ and $(-2i-2) \bmod 9$:

\begin{center}
\begin{tabular}{c|c}
node & links to \\ \hline
0 & 8, 7 \\
1 & 6, 5 \\
2 & 4, 3 \\
3 & 2, 1 \\
4 & 0, 8 \\
5 & 7, 6 \\
6 & 5, 4 \\
7 & 3, 2 \\
8 & 1, 0 \\
\end{tabular}
\end{center}

Average path length is $2.17$, and the best possible is $2.0$, so $\theta \le 1/2.17 \approx 0.46$,
close to the best possible $0.5$.

\end{document}
